% Part: first-order-logic
% Chapter: syntax-and-semantics
% Section: formation-sequences

\documentclass[../../../include/open-logic-section]{subfiles}

\begin{document}

\olfileid{pl}{syn}{fseq}

\olsection{تشکیل دے سلسلے}

!!{formula}s نوں استقرائی تعریف راہیں تعریف کرنا، تے نال دی نال
!!{formula}s دیاں خاصیتاں نوں استقرا راہیں ثابت کرن دی تکنیک، اک سوہنا
تے کارگر طریقہ اے۔ پر !!{formula}s دی بناوٹ نوں ہِٹھوں اُتّے ول، قدم
بہ قدم طریقے نال ویکھنا وی فائدہ مند ہو سکدا اے؛ ایتھے اسیں ایہ کم
\emph{تشکیل دے سلسلے} دے تصور نال کردے آں۔

\begin{defn}[فارمولیاں لئی تشکیل دے سلسلے]
\ollabel{defn:fseq-frm}
علامتاں دی لڑیواں دا محدود سلسلہ $\tuple{!A_0,\dotsc,!A_n}$، جیہڑا
زبان~$\Lang L_0$ توں اے، $!A$ لئی \emph{تشکیل دا سلسلہ} اے جے
$!A \ident !A_n$، تے ہر $i \leq n$ لئی یا تاں $!A_i$ ایٹمی فارمولا
اے، یا ایہو جہے $j,k < i$ موجود نیں کہ اگلیاں وچوں کوئی اک گل پوری
ہُندی اے:
\begin{enumerate}
    \tagitem{prvNot}{$!A_i \ident \lnot !A_j$۔}{}%
    \tagitem{prvAnd}{$!A_i \ident (!A_j \land !A_k)$۔}{}%
    \tagitem{prvOr}{$!A_i \ident (!A_j \lor !A_k)$۔}{}%
    \tagitem{prvIf}{$!A_i \ident (!A_j \lif !A_k)$۔}{}%
    \tagitem{prvIff}{$!A_i \ident (!A_j \liff !A_k)$۔}{}%
\end{enumerate}
\end{defn}

\begin{ex}
\[
\tuple{
    \Obj p_0,
    \Obj p_1,
    (\Obj p_1 \land \Obj p_0),
    \lnot (\Obj p_1 \land \Obj p_0)
}
\]
$\lnot (\Obj p_1 \land \Obj p_0)$ لئی تشکیل دا اک سلسلہ اے، تے ایہ وی
اے:
\[
\tuple{
    \Obj p_0,
    \Obj p_1,
    \Obj p_0,
    (\Obj p_1 \land \Obj p_0),
    (\Obj p_0 \lif \Obj p_1),
    \lnot (\Obj p_1 \land \Obj p_0)
}.
\]
%
دوجی مثال توں ویکھیا جا سکدا اے کہ تشکیل دے سلسلیاں وچ `فالتو مال'
وی آ سکدا اے: اوہ فارمولے جیہڑے وادھو ہون یا بناوٹ وچ کوئی حصہ نہ
پاؤن۔
\end{ex}

\begin{prop}
\ollabel{prop:formed}
ہر !!{formula}~$!A$، جیہڑا~$\Frm[L_0]$ وچ اے، اوہدے لئی تشکیل دا کوئی
سلسلہ موجود اے۔
\end{prop}

\begin{proof}
منّو $!A$ ایٹمی اے۔ تاں سلسلہ~$\tuple{!A}$، $!A$ لئی تشکیل دا سلسلہ
اے۔
%
ہُن منّو کہ $!B$ تے~$!C$ دے تشکیل دے سلسلے ترتیب وار
$\tuple{!B_0,\dotsc,!B_n}$ تے $\tuple{!C_0,\dotsc,!C_m}$ نیں۔
%
\begin{enumerate}
  \tagitem{prvNot}{جے $!A \ident \lnot !B$،
      تاں $\tuple{!B_0,\dotsc,!B_n,\lnot !B_n}$
      $!A$ لئی تشکیل دا سلسلہ اے۔}{}
  \tagitem{prvAnd}{جے $!A \ident (!B \land !C)$،
      تاں $\tuple{!B_0,\dotsc,!B_n,!C_0,\dotsc,!C_m,(!B_n \land !C_m)}$
      $!A$ لئی تشکیل دا سلسلہ اے۔}{}
  \tagitem{prvOr}{جے $!A \ident (!B \lor !C)$،
      تاں $\tuple{!B_0,\dotsc,!B_n,!C_0,\dotsc,!C_m,(!B_n \lor !C_m)}$
      $!A$ لئی تشکیل دا سلسلہ اے۔}{}
  \tagitem{prvIf}{جے $!A \ident (!B \lif !C)$،
      تاں $\tuple{!B_0,\dotsc,!B_n,!C_0,\dotsc,!C_m,(!B_n \lif !C_m)}$
      $!A$ لئی تشکیل دا سلسلہ اے۔}{}
  \tagitem{prvIff}{جے $!A \ident (!B \liff !C)$،
      تاں $\tuple{!B_0,\dotsc,!B_n,!C_0,\dotsc,!C_m,(!B_n \liff !C_m)}$
      $!A$ لئی تشکیل دا سلسلہ اے۔}{}
\end{enumerate}
!!{formula}s اُتے استقرا دے اصول نال، ہر !!{formula} دا کوئی تشکیل دا
سلسلہ اے۔
\end{proof}

اسیں اُلٹا رخ وی ثابت کر سکدے آں۔ ایہ گل اہم اے، کیوں جے ایہ وکھاؤندی
اے کہ فارمولیاں نوں تعریف کرن دے ساڈے دونیں طریقے ہم ارز نیں: دونیں
اکّو نتیجے دیندے نیں۔ ایہدا مطلب ایہ وی اے کہ اسیں تشکیل دے سلسلیاں
دی لمبائی اُتے عام استقرا ورت کے فارمولیاں بارے قضیے ثابت کر سکدے آں۔

\begin{lem}
\ollabel{lem:fseq-init}
منّو کہ $\tuple{!A_0,\dotsc,!A_n}$،~$!A_n$ لئی تشکیل دا سلسلہ اے، تے
$k \leq n$۔ تاں $\tuple{!A_0,\dotsc,!A_k}$،~$!A_k$ لئی تشکیل دا سلسلہ
اے۔
\end{lem}

\begin{proof}
مشق۔
\end{proof}

\begin{thm}
\ollabel{thm:fseq-frm-equiv}
$\Frm[L_0]$، زبان~$\Lang L_0$ وچ علامتاں دیاں اوہناں ساریاں لڑیواں
دا سیٹ اے جنہاں دا کوئی تشکیل دا سلسلہ اے۔
\end{thm}

\begin{proof}
$F$ نوں زبان~$\Lang L_0$ وچ علامتاں دیاں اوہناں ساریاں لڑیواں دا سیٹ
منّو جنہاں دا کوئی تشکیل دا سلسلہ اے۔ اسیں
\olref[pl][syn][fseq]{prop:formed} وچ ویکھ چکے آں کہ
$\Frm[L_0] \subseteq F$، سو ہُن اسیں اُلٹا رخ ثابت کردے آں۔

منّو $!A$ دا تشکیل دا سلسلہ $\tuple{!A_0,\dotsc,!A_n}$ اے۔ اسیں ثابت
کردے آں کہ $!A \in \Frm[L_0]$، تے ایہ کم~$n$ اُتے مضبوط استقرا نال
کردے آں۔ ساڈا
استقرائی مفروضہ اے کہ علامتاں دی ہر اوہ لڑی جیہدے تشکیل دے سلسلے دی
لمبائی $m < n$ اے،~$\Frm[L_0]$ وچ اے۔ تشکیل دے سلسلے دی تعریف موجب
یا تاں $!A_n$ ایٹمی اے، یا ایہو جہے $j,k < n$ لازماً موجود نیں کہ
اگلیاں وچوں کوئی اک صورت اے:
\begin{enumerate}
    \tagitem{prvNot}{$!A_n \ident \lnot !A_j$۔}{}%
    \tagitem{prvAnd}{$!A_n \ident (!A_j \land !A_k)$۔}{}%
    \tagitem{prvOr}{$!A_n \ident (!A_j \lor !A_k)$۔}{}%
    \tagitem{prvIf}{$!A_n \ident (!A_j \lif !A_k)$۔}{}%
    \tagitem{prvIff}{$!A_n \ident (!A_j \liff !A_k)$۔}{}%
\end{enumerate}
ہُن اسیں صورتاں نال دلیل دیندے آں۔ جے $!A_n$ ایٹمی اے تاں
$!A_n \in \Frm[L_0]$۔ ایس دی بجائے منّو کہ $!A \equiv
(!A_j \land !A_k)$۔ \olref[pl][syn][fseq]{lem:fseq-init} موجب،
$\tuple{!A_0,\dotsc,!A_j}$ تے $\tuple{!A_0,\dotsc,!A_k}$ ترتیب وار
$!A_j$ تے~$!A_k$ لئی تشکیل دے سلسلے نیں۔ کیوں جے ایہ~$!A$ دے تشکیل دے
سلسلے دے پورے توں چھوٹے مُڈھلے ذیلی سلسلے نیں، ایس لئی دونیں دی
لمبائی~$n$ توں گھٹ اے۔ سو استقرائی مفروضے موجب $!A_j$ تے~$!A_k$،
$\Frm[L_0]$ وچ نیں، تے ایس لئی !!a{formula} دی تعریف موجب
$(!A_j \land !A_k)$ وی اوہدے وچ اے۔ ہور صورتاں ایسّے برابر دلیل نال
ملدیاں نیں۔
\end{proof}

\end{document}
